\documentclass[11pt]{amsart}
\usepackage{amsmath,amssymb,amsthm}
\usepackage[margin=1.1in]{geometry}
\usepackage{array,longtable,booktabs}
\usepackage{hyperref}
\emergencystretch=2em

\newtheorem{theorem}{Theorem}
\newtheorem{lemma}[theorem]{Lemma}
\newtheorem{corollary}[theorem]{Corollary}
\newtheorem{proposition}[theorem]{Proposition}
\theoremstyle{remark}
\newtheorem{remark}[theorem]{Remark}

\newcommand{\Inf}{\mathrm{Inf}}
\newcommand{\E}{\mathbb{E}}
\newcommand{\MAJ}{\mathrm{MAJ}}
\newcommand{\ii}{\mathrm{i}}

\title{The erasure entropy of a forced coordinate under every bias}
\author{Yitzchak Shmalo}
\address{Einstein Institute of Mathematics, The Hebrew University of Jerusalem, Jerusalem, Israel}
\date{}

\begin{document}

\begin{abstract}
A monotone Boolean function $F_{k,t}$, built from a majority function on $k$ inputs with one hidden coordinate $x$ by
substituting threshold functions of $2t$ fresh variables, has $x$ as the only variable lying both in a shortest
minterm and in a shortest maxterm. The ratio $E_x/\Inf_x$ of the erasure entropy to the influence at $x$ was known to
be unbounded when $k$ and $t$ grow together, and to grow linearly in $k$ for $t=1$. We show that for every fixed
$t\ge1$ it equals $(1+o(1))\log_2(1/\Inf_x)=(1+o(1))(k-1)\log_2(1/P_t)$, where $P_t$ is the probability that a binomial
variable with parameters $2t$ and $\frac12$ is at least $t$. More generally the same holds for the underlying majority
function under every rational product bias $P\in(\frac12,1)$, the phase of the relevant coefficients advancing by
$\arccos\frac{2P-1}{2P}$, an irrational multiple of $\pi$. At $P=\frac12$ that angle is $\frac\pi2$, and the behaviour
changes: there the coefficients are values of Krawtchouk polynomials at the middle of their range, half of them
vanish, and the ratio is $\frac12\log_2(1/\Inf_x)-\frac12\log_2\log_2(1/\Inf_x)+c_0+\frac12+O(1/\log_2(1/\Inf_x))$ with
$c_0=\frac34\log_2e-\frac12\log_2\frac\pi2$, which we derive from an exact single-sum formula.
\end{abstract}

\subjclass[2020]{06E30, 42C10, 94A17}
\keywords{Fourier entropy-influence conjecture, Boolean functions, Fourier entropy, total influence}

\maketitle

\section{Introduction}\label{sec:intro}

For $f\colon\{0,1\}^n\to\{-1,1\}$ write $f=\sum_S\hat f(S)\chi_S$ and let $U$ be a random subset with
$\Pr[U=S]=\hat f(S)^2$. The Fourier entropy-influence conjecture of Friedgut and Kalai~\cite{FK96} asks for an absolute
constant $C_{\mathrm{FEI}}$ with $H(U)\le C_{\mathrm{FEI}}\,\E|U|$. It is known for symmetric functions and read-once decision trees~\cite{OWZ} and is
preserved under suitable compositions~\cite{OT13}. Weaker forms and bounds with extra factors are
in~\cite{CKLS,WWW,KKLMS,Han}. Li and Han~\cite{LH25} make the constant in front of $\E|U|$ in Han's bound
sharp. The conjecture was proved for every linear threshold function in~\cite{ShmLTF}, and the min-entropy version is
known for read-$k$ DNFs with $k$ fixed~\cite{Shalev,ACKSW}. The largest known lower bound on the best constant is
$6.5218$~\cite{Num26}. Earlier bounds are due to O'Donnell and Tan~\cite{OT13}, Hod~\cite{Hod} and Mosaic
Intelligence~\cite{MI26a,MI26b}. The erasure entropy of $f$ at a
coordinate $x$ is $E_x(f)=H(\mathbf 1[x\in U]\mid U\setminus\{x\})$, the conditional entropy of one coordinate given
the others in the sense of Verd\'u and Weissman~\cite{VW08}, and with $\Inf_x(f)=\Pr[x\in U]$ the conjecture holds if
and only if $\sup_f\min_{x:\Inf_x(f)>0}E_x(f)/\Inf_x(f)$ is finite~\cite{Shm1}, which builds on the
reduction of Gonz\'alez, MacManus and Pereyra~\cite{GMP} to a statement about one coordinate. Since conditioning does not increase
entropy,
\begin{equation}\label{eq:upper}
\frac{E_x(f)}{\Inf_x(f)}\le\frac{h_2(\Inf_x(f))}{\Inf_x(f)}\le\log_2\frac{e}{\Inf_x(f)}
\end{equation}
for every $f$ and $x$ with $\Inf_x(f)>0$.

The question is which coordinate to take. In~\cite{Shm8} several selection rules are shown to fail on monotone
functions, among them the rule that takes, among the variables lying both in a shortest minterm and in a shortest
maxterm, one of smallest ratio. Fix an odd $k=2m+1$, let $y^0\in\{0,1\}^k$ have $y^0_j=1$ for $j\le m$ and
$y^0_j=0$ for $j>m$, let $y^1$ differ from $y^0$ only in $y^1_{m+1}=1$, let $M$ be $\MAJ_k$ with the point $y^1$ removed from its true set, and
\[
\tilde f_k(y,x)=M(y)\vee\bigl(x\wedge[y\ge y^0]\bigr).
\]
For $t\ge1$ the function $F_{k,t}$ of~\cite[Section~5]{Shm8} replaces $y_j$ by the threshold function ``at least $t$
of $2t$'' of a block of fresh variables for $j\le m$, and by ``at least $t+1$ of $2t$'' for $j>m$. It is monotone, and
$x$ is the only variable lying both in a shortest minterm and in a shortest maxterm. Moreover~\cite[proof of
Proposition~12]{Shm8} $E_x(F_{k,t})$ and $\Inf_x(F_{k,t})$ equal the erasure entropy and the influence of $x$ in the
expansion of $\tilde f_k$ in the $\mu$-biased Fourier basis~\cite[Chapter~8]{OD14} of the product measure with
$\Pr[y_j=1]=P_t$ for $j\le m$ and $1-P_t$ for $j>m$, where
\[
P_t=\Pr\bigl[\mathrm{Bin}(2t,\tfrac12)\ge t\bigr]=\frac12+\binom{2t}{t}2^{-2t-1};
\]
the biased form of the conjecture is in~\cite{KMS12}, and a sharp lower bound on the biased Fourier entropy by a
sum of coordinate-wise terms is in~\cite{Cha25}. For $t=1$, $P_1=\frac34$ and the ratio at $x$ is
$(1+o(1))(k-1)\log_2\frac43$~\cite{Shm9}. Here we treat every fixed $t\ge1$, and more generally every rational bias.

For $P\in(0,1)$ write $E_x^P(\tilde f_k)$ and $\Inf_x^P(\tilde f_k)$ for the erasure entropy and the influence of $x$
for the law that gives each set $S$ the square of the coefficient of $\tilde f_k$ at $S$, in $\pm1$ values, in the orthonormal basis of
the product measure with $\Pr[y_j=1]=P$ for $j\le m$, $\Pr[y_j=1]=1-P$ for $j>m$, and $x$ uniform. At $P=\frac12$
these are $E_x(\tilde f_k)$ and $\Inf_x(\tilde f_k)$. Let $\alpha_P=\arccos\frac{2P-1}{2P}$.

\begin{theorem}\label{thm:bias}
Let $P\in(\frac12,1)$ with $\alpha_P/\pi$ irrational. Then $\Inf_x^P(\tilde f_k)=P^{k-1}$ and, as $k\to\infty$ through
odd values,
\[
\frac{E_x^P(\tilde f_k)}{\Inf_x^P(\tilde f_k)}=(1+o(1))\log_2\frac1{\Inf_x^P(\tilde f_k)}=(1+o(1))\,(k-1)\log_2\frac1P .
\]
The condition holds for every rational $P\in(\frac12,1)$.
\end{theorem}

\begin{corollary}\label{cor:t}
For every fixed $t\ge1$, $\Inf_x(F_{k,t})=P_t^{k-1}$ and $E_x(F_{k,t})/\Inf_x(F_{k,t})=(1+o(1))(k-1)\log_2(1/P_t)$ as
$k\to\infty$. In particular the minterm-maxterm rule of~\cite{Shm8} has unbounded ratio on the functions $F_{k,t}$
with $t$ fixed.
\end{corollary}

The corollary is Theorem~\ref{thm:bias} with $P=P_t$, which is rational, together with the reduction quoted above.
The case $P=\frac34$ is the theorem of~\cite{Shm9}, and the proof below follows that paper; what makes it work for
every $P$ are three identities on the unit circle, \eqref{eq:ident} below.

At $P=\frac12$ the rotation angle $\alpha_{1/2}$ is $\frac\pi2$ and the argument breaks down for a reason that can be
made exact. Write $\Psi(r)=r^2\log_2(1+r^{-2})+\log_2(1+r^2)$ for $r>0$ and $\Psi(0)=0$.

\begin{theorem}\label{thm:half}
Let $k=2m+1$. Then $\Inf_x(\tilde f_k)=2^{-2m}$ and
\[
\frac{E_x(\tilde f_k)}{\Inf_x(\tilde f_k)}=4^{-m}\sum_{b=0}^{m-1}\binom{2m}{2b+1}\,
\Psi\Bigl(\binom mb\binom{2m}m\Big/\binom{2m}{2b}\Bigr).
\]
Moreover $\bigl|E_x(\tilde f_k)/\Inf_x(\tilde f_k)-m+\frac12\log_2m\bigr|\le C$ for an absolute constant $C$, and
\[
\frac{E_x(\tilde f_k)}{\Inf_x(\tilde f_k)}=m-\frac12\log_2m+c_0+O\Bigl(\frac1m\Bigr),\qquad
c_0=\frac34\log_2e-\frac12\log_2\frac\pi2=0.75627\dots,
\]
so that, as $\log_2(1/\Inf_x(\tilde f_k))=2m$, the ratio is $\frac12\log_2\frac1{\Inf_x}-\frac12\log_2\log_2\frac1{\Inf_x}+c_0+\frac12+O(1/m)$.
\end{theorem}

In~\cite[Proposition~11]{Shm8} the ratio of $\tilde f_k$ was shown to be at least $(\frac18-o(1))(k-1)$; the
theorem gives the constant $\frac12$. The factor $\frac12$ has a direct meaning: on exactly half of the spectral weight
of the derivative in $x$ the average over $x$ has a vanishing coefficient, so there $x\in U$ is determined by
$U\setminus\{x\}$ and contributes no entropy. As $t\to\infty$ with $k$ fixed the ratio of $F_{k,t}$ tends to that of
$\tilde f_k$~\cite[Proposition~12]{Shm8}, while for each fixed $t$ the ratio is asymptotic to the bound
\eqref{eq:upper}; the two limits do not commute.

\section{The ratio as a double sum}\label{sec:sum}

All logarithms are to base $2$ unless written $\ln$. For $a,b\ge0$ let $\Phi(a,b)=a\log\frac{a+b}{a}+b\log\frac{a+b}{b}$,
with $\Phi(a,0)=\Phi(0,b)=0$, so that $\Psi(r)=\Phi(r^2,1)$. Then $\Psi\ge0$, $\Psi(r)\ge\log r^2$ for $r\ge1$, and
$\Psi(r)\le\log r^2+1+\log e$ for $r\ge1$.

Fix $P\in[\frac12,1)$ and put $q=1-P$ and $\sigma=\sqrt{Pq}$. In the basis $\prod_{j\in S}\phi_j(y_j)$, with
$\phi_j(y_j)=(y_j-\mu_j)/\sigma$, $\mu_j=P$ for $j\le m$ and $\mu_j=q$ for $j>m$, and with $+1$ for true,
$\tilde f_k=\bar A+(2x-1)D$, where
\[
D=\mathbf 1[y=y^0]+\mathbf 1[y=y^1],\qquad \bar A=2N-1+\mathbf 1[y=y^0]-\mathbf 1[y=y^1],\qquad
N=\mathbf 1[|y|\ge m+1].
\]
For $S\subseteq[k]$ let $u=|S\cap\{1,\dots,m\}|$ and $v=|S\cap\{m+2,\dots,k\}|$, and put
\begin{gather*}
X(z)=q+Pz,\qquad Y(z)=P+qz,\\
\Pi_{uv}(z)=(z-1)^{u+v}X(z)^{m-u}Y(z)^{m+1-v},\qquad \tau_{uv}=\sum_{\nu\ge m+1}[z^\nu]\Pi_{uv}(z).
\end{gather*}

\begin{lemma}\label{lem:coef}
If $m+1\in S$ then $\hat D(S)=0$. If $m+1\notin S$ then $\hat D(S)=(-1)^v\sigma^{u+v}P^{2m-u-v}$ and
$\hat{\bar A}(S)=2\sigma^{u+v}\tau_{uv}-\mathbf 1[S=\emptyset]+(P-q)\hat D(S)$.
\end{lemma}

\begin{proof}
This is~\cite[Lemma~1]{Shm9} with $P$ in place of $\frac34$. For a point $y'$,
$\widehat{\mathbf 1[y=y']}(S)=\Pr[y=y']\prod_{j\in S}\phi_j(y'_j)$, and $\Pr[y=y^0]=P^{2m+1}$, $\Pr[y=y^1]=P^{2m}q$. If
$m+1\in S$ the two terms of $\hat D(S)$ are $P^{2m+1}\phi_{m+1}(0)=-P^{2m}\sigma$ and $P^{2m}q\,\phi_{m+1}(1)=P^{2m}\sigma$
times the same product, and cancel. Otherwise both points give the product $(-1)^v(q/P)^{(u+v)/2}$, since
$\phi_j(1)=\sqrt{q/P}$ for $j\le m$ and $\phi_j(0)=-\sqrt{q/P}$ for $j\ge m+2$; this gives $\hat D(S)$, and the
difference of the two indicators has coefficient $(P-q)\hat D(S)$. Finally $\E[z^{y_j}\phi_j(y_j)]=\sigma(z-1)$ and
$\E[z^{y_j}]=1-\mu_j+\mu_jz$, so $\E[z^{|y|}\prod_{j\in S}\phi_j(y_j)]=\sigma^{u+v}\Pi_{uv}(z)$ when $m+1\notin S$, and
$\hat N(S)$ is the sum of its coefficients at $z^\nu$, $\nu\ge m+1$.
\end{proof}

Let $\beta_u=\binom mu q^uP^{m-u}$ and, for $0\le u,v\le m$,
\[
r_{uv}=\frac{\bigl|2\tau_{uv}-\mathbf 1[u=v=0]+(P-q)(-1)^vP^{2m-u-v}\bigr|}{P^{2m-u-v}}.
\]

\begin{lemma}\label{lem:sum}
$\Inf_x^P(\tilde f_k)=P^{2m}$ and $E_x^P(\tilde f_k)/\Inf_x^P(\tilde f_k)=\sum_{u,v=0}^m\beta_u\beta_v\,\Psi(r_{uv})$.
\end{lemma}

\begin{proof}
As in~\cite[Lemma~2]{Shm9}: $E_x=\sum_S\Phi(\hat{\bar A}(S)^2,\hat D(S)^2)$ and $\Inf_x=\sum_S\hat D(S)^2$, the terms
with $m+1\in S$ vanish, there are $\binom mu\binom mv$ sets with given $(u,v)$ and $m+1\notin S$,
$\sum_{u,v}\binom mu\binom mv(Pq)^{u+v}P^{2(2m-u-v)}=(Pq+P^2)^{2m}=P^{2m}$, and
$\Phi(\hat{\bar A}^2,\hat D^2)=\hat D^2\Psi(|\hat{\bar A}|/|\hat D|)$ with $|\hat{\bar A}(S)|/|\hat D(S)|=r_{uv}$.
\end{proof}

If $u+v\ge1$ then $Q_{uv}=\Pi_{uv}/(z-1)$ is a polynomial and $\tau_{uv}=[z^m]Q_{uv}$, since $[z^\nu]\Pi_{uv}=[z^{\nu-1}]Q_{uv}-[z^\nu]Q_{uv}$ and the
sum over $\nu\ge m+1$ telescopes~\cite[Lemma~3]{Shm9}.

\section{The unbiased case}\label{sec:half}

Here $P=q=\frac12$, so $X=Y=(1+z)/2$, the term $(P-q)\hat D$ is absent, and $\beta_u=\binom mu2^{-m}$. For
$0\le a\le 2m$ let
\[
\kappa_a=[z^m](1-z)^a(1+z)^{2m-a},
\]
the value at $a$ of the Krawtchouk polynomial of degree $m$ for length $2m$~\cite[Chapter~5, \S7]{MS77}.

\begin{lemma}\label{lem:kraw}
$\kappa_a=0$ for odd $a$, and $\kappa_{2b}=(-1)^b\binom mb\binom{2m}m\big/\binom{2m}{2b}$.
\end{lemma}

\begin{proof}
For a variable $\xi$,
\[
\sum_{a=0}^{2m}\binom{2m}a\kappa_a\xi^a=[z^m]\bigl((1+z)+\xi(1-z)\bigr)^{2m}=[z^m]\bigl((1+\xi)+(1-\xi)z\bigr)^{2m}
=\binom{2m}m(1-\xi^2)^m .
\]
Comparing coefficients of $\xi^a$ gives the claim.
\end{proof}

\begin{proof}[Proof of Theorem~\ref{thm:half}]
By Lemma~\ref{lem:sum}, $\Inf_x=2^{-2m}$. For $u+v\ge1$ put $a=u+v-1$; then
\[
\frac{2\tau_{uv}}{P^{2m-u-v}}=2^{2m-u-v+1}\,[z^m](z-1)^a\Bigl(\frac{1+z}2\Bigr)^{2m-a}=(-1)^a\kappa_a ,
\]
so $r_{uv}=|\kappa_{u+v-1}|$. For $u=v=0$, $\tau_{00}=\Pr[|y|\ge m+1]=\frac12$ and $r_{00}=0$. The weight of the pairs
with $u+v=s$ is $\sum_{u+v=s}\beta_u\beta_v=\binom{2m}s4^{-m}$, and by Lemma~\ref{lem:kraw} only odd $s=2b+1$ give
$r\ne0$, with $r=\binom mb\binom{2m}m/\binom{2m}{2b}$. This is the formula. Since $\binom{2m}{2b}\le\binom{2m}m$, these
values satisfy $r\ge\binom mb\ge1$.

For the asymptotics write $p_b=\binom{2m}{2b+1}4^{-m}$ and $\delta_b=(2b-m)^2/m$. The odd binomial coefficients of
order $2m$ sum to $2^{2m-1}$, so $\sum_bp_b=\frac12$, and if $s$ is binomial with parameters $2m$ and $\frac12$ then
$\sum_bp_b\delta_b\le\E[(s-1-m)^2]/m=\frac12+\frac1m$. We use $h_2(\frac12+\zeta)\ge1-4\zeta^2$, $|\zeta|\le\frac12$, for the binary entropy $h_2$
and, for $1\le j\le n-1$,
\begin{equation}\label{eq:binom}
\frac{2^{nh_2(j/n)}}{\sqrt{8j(n-j)/n}}\le\binom nj\le\frac{2^{nh_2(j/n)}}{\sqrt{2\pi j(n-j)/n}}
\end{equation}
\cite[Chapter~10, Lemma~7]{MS77}, together with $\binom{2m}m\le4^m/\sqrt{\pi m}$.

Lower bound. By~\eqref{eq:binom} with $n=m$ and $8b(m-b)/m\le2m$, $\log\binom mb\ge m\,h_2(b/m)-\frac12\log(2m)\ge
m-\delta_b-\frac12\log m-\frac12$ for $1\le b\le m-1$, and this also holds for $b=0$. As $\Psi(r)\ge\log r^2$,
\[
\frac{E_x}{\Inf_x}\ge2\sum_bp_b\log\binom mb\ge m-\tfrac12\log m-\tfrac32-\tfrac2m .
\]

Upper bound. For $1\le b\le m-1$, \eqref{eq:binom} applied to $\binom mb$ from above and to $\binom{2m}{2b}$ from below
gives
\begin{align*}
\log r=\log\binom mb+\log\binom{2m}m-\log\binom{2m}{2b}
&\le 2m-m\,h_2(b/m)-\tfrac12\log m+\tfrac32-\log\pi\\
&\le m+\delta_b-\tfrac12\log m+1 ,
\end{align*}
the factors $b(m-b)/m$ cancelling. For $b=0$, $r=\binom{2m}m\le4^m$ and $p_0=2m4^{-m}$. With
$\Psi(r)\le\log r^2+1+\log e$,
\[
\frac{E_x}{\Inf_x}\le\sum_{b\ge1}p_b\bigl(2m+2\delta_b-\log m+3+\log e\bigr)+p_0(4m+3)
\le m-\tfrac12\log m+C .
\]
Finally $\log_2(1/\Inf_x)=2m$, and $\frac12\log m=\frac12\log(2m)-\frac12$.

The constant. Let $m\ge2$, write $d_b=2b-m$, so that $\delta_b=d_b^2/m$, and let $C$ denote absolute constants.
First, $r$ is large for every $b$: since $j\mapsto\binom mj$ is log-concave, each term of
$\binom{2m}{2b}=\sum_j\binom mj\binom m{2b-j}$ is at most $\binom mb^2$, so $\binom{2m}{2b}\le(m+1)\binom mb^2$ and,
with $\binom{2m}m\ge4^m/(2\sqrt m)$ and $\binom mb\le2^m$,
\[
r=\binom mb\binom{2m}m\Big/\binom{2m}{2b}\ge\frac{\binom{2m}m}{(m+1)\binom mb}\ge\frac{2^m}{2(m+1)\sqrt m}.
\]
For $r\ge1$, $\Psi(r)-\log r^2-\log e=\log e\,\bigl(r^2\ln(1+r^{-2})-1+\ln(1+r^{-2})\bigr)$ lies between
$-\frac12r^{-2}\log e$ and $r^{-2}\log e$. Hence
\begin{equation}\label{eq:psilog}
\frac{E_x}{\Inf_x}=\frac12\log e+\sum_bp_b\log r^2+O(m^3 4^{-m}).
\end{equation}
Next let $B=\{b:|d_b|\le m/4\}$. For $b\in B$ we use $\ln\binom nj=n\,h_e(j/n)-\frac12\ln\bigl(2\pi j(n-j)/n\bigr)
+O(1/n)$ for $\frac38\le j/n\le\frac58$, which follows from Robbins' form of Stirling's formula~\cite{Rob55}, where
$h_e(\frac12+\zeta)=\ln2-2\zeta^2+O(\zeta^4)$ is the binary entropy in nats. With $\zeta=d_b/(2m)$ for both
$\binom mb$ and $\binom{2m}{2b}$, and $\ln(1-4\zeta^2)=O(\zeta^2)$,
\[
\ln\binom mb=m\ln2-\frac{d_b^2}{2m}-\frac12\ln\frac{\pi m}2+\Delta_1,\qquad
\ln\binom{2m}{2b}=2m\ln2-\frac{d_b^2}m-\frac12\ln(\pi m)+\Delta_2,
\]
and $\ln\binom{2m}m=2m\ln2-\frac12\ln(\pi m)+\Delta_3$, with
$|\Delta_i|\le C(d_b^4/m^3+d_b^2/m^2+1/m)$. Therefore, for $b\in B$,
\[
\log r^2=2m-\log\frac{\pi m}2+\delta_b\log e+\Delta_b,\qquad|\Delta_b|\le C\Bigl(\frac{d_b^4}{m^3}
+\frac{d_b^2}{m^2}+\frac1m\Bigr).
\]
For $b\notin B$ both $\log r^2$ and $2m-\log\frac{\pi m}2+\delta_b\log e$ lie in $[-\log m,6m]$, because
$1\le r\le\binom mb\binom{2m}m\le8^m$ and $\delta_b\le m$. If $s$ is binomial with parameters $2m$ and $\frac12$, then
$p_b=\Pr[s=2b+1]$ and $d_b=s-1-m$, and for $m\ge4$ Hoeffding's inequality~\cite[Theorem~2]{Hoe63} gives
$\sum_{b\notin B}p_b\le\Pr[|s-m|>\frac m4-1]\le2e^{-(m-4)^2/(16m)}$. Moreover $\E[(s-m)^4]=O(m^2)$, so $\sum_bp_b|\Delta_b|=O(1/m)$ over $b\in B$. Finally,
$\sum_s(-1)^s\binom{2m}ss^j=0$ for $0\le j<2m$, so the odd and the even values of
$s$ carry equal parts of $\E[(s-1-m)^2]=\frac m2+1$, and
\[
\sum_bp_b\delta_b=\frac1m\cdot\frac12\Bigl(\frac m2+1\Bigr)=\frac14+\frac1{2m}.
\]
Putting these together with $\sum_bp_b=\frac12$,
\[
\sum_bp_b\log r^2=\frac12\Bigl(2m-\log\frac{\pi m}2\Bigr)+\frac14\log e+O\Bigl(\frac1m\Bigr),
\]
and by~\eqref{eq:psilog} $E_x/\Inf_x=m-\frac12\log m+\frac34\log e-\frac12\log\frac\pi2+O(1/m)$.
\end{proof}

\begin{remark}
The proof shows where the weight sits: $r=0$ exactly when $u+v$ is even, which is half of the weight, and on the other
half $\log r^2$ is $2m-\log m+O(1)$ for typical $b$. The constant $c_0$ collects three contributions: the
Gaussian shape of the three binomial coefficients, which gives $-\frac12\log\frac\pi2$; their curvature, which makes
$\log r^2$ grow by $\delta_b\log e$ away from the centre and contributes $\frac14\log e$; and the gap
$\Psi(r)-\log r^2\to\log e$, which contributes $\frac12\log e$.
\end{remark}

\begin{proposition}\label{prop:next}
In the setting of Theorem~\ref{thm:half},
\[
\frac{E_x(\tilde f_k)}{\Inf_x(\tilde f_k)}=m-\frac12\log_2m+c_0+\frac{5\log_2e}{16\,m}+O\Bigl(\frac1{m^2}\Bigr).
\]
\end{proposition}

\begin{proof}
By~\eqref{eq:psilog} and the tail bound in the proof of Theorem~\ref{thm:half}, it suffices to expand
$\ln r$ for $|d_b|\le m/4$ and average. Stirling's series with its first correction,
$\ln n!=n\ln n-n+\frac12\ln(2\pi n)+\frac1{12n}+O(n^{-3})$, gives for $\frac38\le j/n\le\frac58$
\[
\ln\binom nj=n\,h_e(j/n)-\frac12\ln\frac{2\pi j(n-j)}n+\frac1{12}\Bigl(\frac1n-\frac1j-\frac1{n-j}\Bigr)+O(n^{-3}).
\]
Apply this to $\binom mb$, $\binom{2m}m$ and $\binom{2m}{2b}$ with $\zeta=d_b/(2m)$ in the first and third. The
entropy terms give $m\ln2-m\bigl(h_e(\frac12+\zeta)-\ln2\bigr)=m\ln2+2m\zeta^2+\frac43m\zeta^4+O(m\zeta^6)$, because
$h_e(\frac12+\zeta)=\ln2-2\zeta^2-\frac43\zeta^4+O(\zeta^6)$. In the logarithmic terms the factors $1-4\zeta^2$ of
$\binom mb$ and $\binom{2m}{2b}$ cancel, leaving $-\frac12\ln\frac{\pi m}2$. The brackets of the correction terms
are $-\frac3m$, $-\frac3{2m}$ and $-\frac3{2m}$ up to $O(d_b^2/m^3)$; they enter with signs $+,+,-$, so together the
corrections give $-\frac1{4m}$. Hence
\[
\ln r=m\ln2+\frac{d_b^2}{2m}+\frac{d_b^4}{12m^3}-\frac12\ln\frac{\pi m}2-\frac1{4m}
+O\Bigl(\frac{d_b^2}{m^3}+\frac{d_b^6}{m^5}+\frac1{m^3}\Bigr).
\]
With $s$ binomial with parameters $2m$ and $\frac12$ as before, the alternating-sum identity gives
$\sum_bp_bd_b^2=\frac12\E[(s-1-m)^2]=\frac m4+\frac12$ and $\sum_bp_bd_b^4=\frac12\E[(s-1-m)^4]=\frac{3m^2}8+O(m)$, the
latter because the central fourth moment of $s$ is $\frac{m(3m-1)}4$. The averages of the error terms are
$O(1/m^2)$. Therefore
\[
\sum_bp_b\ln r=\frac m2\ln2-\frac14\ln\frac{\pi m}2+\frac18+\Bigl(\frac14+\frac1{32}-\frac18\Bigr)\frac1m
+O\Bigl(\frac1{m^2}\Bigr),
\]
and $E_x/\Inf_x=\frac12\log e+\frac2{\ln2}\sum_bp_b\ln r+O(m^3 4^{-m})$ gives the claim, the coefficient of $1/m$
being $\frac2{\ln 2}\cdot\frac5{32}=\frac5{16}\log e$.
\end{proof}

\begin{remark}\label{rem:near}
Near $P=\frac12$ the two regimes meet. Write $P=\frac12+\eta$. At $\eta=0$,
$\partial_\eta \Pi_{uv}=(v-u-1)(z-1)^{u+v+1}\bigl(\frac{1+z}2\bigr)^{2m-u-v}$, so the terms with $u+v=2b+1$ keep
$r_{uv}=|\kappa_{2b}|+O(\eta)$, while those with $u+v=2b\ge2$, which vanish at $\eta=0$, have
\[
r_{uv}=2\eta\bigl|(v-u-1)\kappa_{2b}+(-1)^v\bigr|+O(\eta^2),
\]
where $v-u-1$ is odd, hence never $0$. When $\eta$ is small and $\eta2^m$ is large, all these $r_{uv}$ are
large, $\Psi(r_{uv})=\log r_{uv}^2+\log e+o(1)$, and to first order the sum of Lemma~\ref{lem:sum} is
$(k-1)\log\frac1P-\log\frac1\eta-\frac12\log m+c_1(m)$ with $c_1(m)$ independent of $\eta$; the leading term
$(k-1)\log\frac1P$ differs from $2m$ by $O(m\eta)$. The weight of the pairs with $u+v$ even is
$\frac12-4^{-m}$, $v-u$ is close to normal with variance $\frac m2$ and nearly independent of $u+v$, and
$\E\ln Z^2=-\gamma-\ln2$ for a standard normal $Z$, where $\gamma$ is Euler's constant. With the computation in the
proof of Theorem~\ref{thm:half} for the factors $\kappa_{2b}$, this gives
\[
c_1(m)=\log\frac2\pi+\frac{3-\gamma}2\log e+\frac2{\sqrt{\pi m}}+O\Bigl(\frac1m\Bigr),\qquad
\log\frac2\pi+\frac{3-\gamma}2\log e=1.096173\dots,
\]
the term $2/\sqrt{\pi m}$ coming from the lattice of odd values of $v-u-1$: for even $K$ the sum of $\ln j^2$ over the odd $j$ with
$|j|<K$ exceeds $2K\ln K-2K$ by $2\ln2+o(1)$, and the density of $v-u$ at $0$ is about
$1/\sqrt{\pi m}$. We have not made the error terms rigorous. Evaluated
exactly in the first-order form, $c_1(m)$ is $1.423333$, $1.353700$, $1.313931$, $1.225405$ and $1.156825$ at $m=20$,
$30$, $40$, $100$ and $400$, and $(c_1(m)-1.096173)\sqrt m$ is $1.463$, $1.411$, $1.377$, $1.292$ and $1.213$, approaching
$2/\sqrt\pi=1.128$ at the rate $1/\sqrt m$. For the exact sum of Lemma~\ref{lem:sum} at $k=61$ and
$\eta=10^{-4}$, $2^{-15}$ and $10^{-6}$, the ratio minus $(k-1)\log\frac1P-\log\frac1\eta-\frac12\log m$
is $1.3540$, $1.3537$ and $1.3537$, against $c_1(30)=1.353700$, and at $k=81$ and $\eta=10^{-4}$, $10^{-6}$
and $2^{-20}$ it is $1.3142$, $1.3139$ and $1.3139$, against $c_1(40)=1.313931$. Once $\eta2^m$ is of order $1$ the first-order form fails, and the ratio
moves to the value of Theorem~\ref{thm:half}.
\end{remark}

\section{The tail on the unit circle}\label{sec:circle}

From now on $P\in(\frac12,1)$. For $z=e^{\ii\theta}$ and $w=\cos\theta$,
\[
|z-1|^2=2-2w,\qquad |X(z)|^2=|Y(z)|^2=R(w):=P^2+q^2+2Pqw,
\]
and $X(z)=z\,\overline{Y(z)}$. Put $n_1=u+v-1$ and $n_2=2m+1-u-v$, so $n_1+n_2=2m$. Then
$|Q_{uv}(e^{\ii\theta})|=e^{\rho(\cos\theta)}$ with
\[
\rho(w)=\frac{n_1}2\ln(2-2w)+\frac{n_2}2\ln R(w),\qquad
\rho''(w)=-\frac{n_1}{2(1-w)^2}-\frac{2P^2q^2n_2}{R(w)^2}\le-2P^2q^2n_2 ,
\]
using $R\le1$. So $\rho$ is strictly concave on $[-1,1)$, with maximum at
\[
w_0=\frac{2Pq\,n_2-(P^2+q^2)\,n_1}{4Pq\,m}
\]
when this lies in $(-1,1)$. A short computation gives $2-2w_0=n_1/(2Pqm)$ and $R(w_0)=n_2/(2m)$, hence
\begin{equation}\label{eq:peak}
\rho(w_0)=\frac{n_1}2\ln\frac{n_1}{2Pqm}+\frac{n_2}2\ln\frac{n_2}{2m}.
\end{equation}
The law $\beta_u$ has mean $mq$ and variance $mPq$. For $\lambda\ge1$ let $W_\lambda$ be the set of $(u,v)$ with
$|u-mq|\le\lambda\sqrt m$ and $|v-mq|\le\lambda\sqrt m$. Below, $c,C$ are positive constants that may depend on $P$
and $\lambda$ but not on $m$ or on $(u,v)\in W_\lambda$, and $m$ is large enough for each statement. On $W_\lambda$,
$n_1=2mq+O(\sqrt m)$ and $w_0=w^*+O(m^{-1/2})$, where
\[
w^*=\frac{2P-1}{2P}=\cos\alpha_P\in(0,\tfrac12).
\]
Put $\theta_0=\arccos w_0$, $z_0=e^{\ii\theta_0}$, $\theta^*=\alpha_P$ and $z^*=e^{\ii\theta^*}$. Since
$1-w^*=1/(2P)$, the real part of $(z^*-1)(q+P\bar z^*)$ is $(P-q)(1-w^*)=w^*$ and its imaginary part is $\sin\theta^*$, and
$R(w^*)=P^2+q^2+q(P-q)=P$; so
\begin{equation}\label{eq:ident}
(z^*-1)\,\overline{X(z^*)}=z^*,\qquad |X(z^*)|^2=P,\qquad Y(z^*)=z^*\,\overline{X(z^*)} .
\end{equation}

\begin{lemma}\label{lem:exponent}
On $W_\lambda$, $\rho(w_0)+(n_2-1)\ln\frac1P\ge(m-1)\ln\frac1P$.
\end{lemma}

\begin{proof}
Regard the left side, by~\eqref{eq:peak}, as a function $G(n_1)$ of a real $n_1\in(0,2m)$ with $n_2=2m-n_1$. Then
$G'(n_1)=\frac12\ln\frac{n_1}{Pq\,n_2}+\ln P$, which vanishes at $n_1=2mq$, and $G''=\frac1{2n_1}+\frac1{2n_2}>0$.
So $G\ge G(2mq)=mq\ln\frac1P+mP\ln P+(2mP-1)\ln\frac1P=(m-1)\ln\frac1P$.
\end{proof}

Fix $(u,v)\in W_\lambda$ and let $\Lambda(\theta)=\ln Q_{uv}(e^{\ii\theta})-\ii m\theta$ for complex $\theta$ in a closed
disc centred at $\theta^*$ and so small that $e^{\ii\theta}$ stays away from $1$, $-q/P$ and $-P/q$,
the zeros of $z-1$, $X$ and $Y$; the logarithm is the branch continuous in this disc. Then $|\Lambda''|,|\Lambda'''|\le Cm$
on the disc. For real $\theta$, $\mathrm{Re}\,\Lambda(\theta)=\rho(\cos\theta)$, so $\mathrm{Re}\,\Lambda'(\theta_0)=0$ and
$\mathrm{Re}\,\Lambda''(\theta_0)=\rho''(w_0)\sin^2\theta_0\le-cm$. Put
\[
\omega=\mathrm{Im}\,\Lambda'(\theta_0),\qquad \Omega=-\Lambda''(\theta_0),\qquad
T_{uv}=\frac{1}{\sqrt{2\pi\Omega}}\exp\Bigl(\Lambda(\theta_0)-\frac{\omega^2}{2\Omega}\Bigr).
\]

\begin{lemma}\label{lem:laplace}
On $W_\lambda$: $|\omega|\le C\sqrt m$, $|T_{uv}|\ge c\,m^{-1/2}e^{\rho(w_0)}$, and
$\tau_{uv}=2\,\mathrm{Re}\,T_{uv}+O(|T_{uv}|\,m^{-1/5})$.
\end{lemma}

\begin{proof}
The function $\Lambda'(\theta)/m$ differs by $O(1/m)$ from
\[
\ii z\Bigl(\frac{(u+v)/m}{z-1}+\frac{(1-u/m)P}{X(z)}+\frac{(1-v/m)q}{Y(z)}\Bigr)-\ii,\qquad z=e^{\ii\theta},
\]
a smooth function of $(u/m,v/m,\theta)$. It vanishes at $u/m=v/m=q$, $\theta=\theta^*$: by~\eqref{eq:ident},
$1/(z^*-1)=\overline{X^*}/z^*$, $1/X^*=\overline{X^*}/P$ and $1/Y^*=X^*/(Pz^*)$, where $X^*=X(z^*)$, so
\[
z^*\Bigl(\frac{2q}{z^*-1}+\frac{P^2}{X^*}+\frac{Pq}{Y^*}\Bigr)=2q\overline{X^*}+Pz^*\overline{X^*}+qX^*
=3q^2+P^2+4Pq\,w^*=1 .
\]
The rest is the proof of~\cite[Lemma~5]{Shm9}. Since $u/m-q$, $v/m-q$ and $\theta_0-\theta^*$ are $O(m^{-1/2})$,
$|\Lambda'(\theta_0)|\le C\sqrt m$, and $\mathrm{Re}(\omega^2/\Omega)=\omega^2\mathrm{Re}\,\Omega/|\Omega|^2\le C$, which gives the lower bound
on $|T_{uv}|$. By Cauchy's formula~\cite[Chapter~VIII]{FS09}, $\tau_{uv}=\frac1{2\pi}\int_{-\pi}^\pi e^{\Lambda(\theta)}d\theta$.
On $J_+=\{|\theta-\theta_0|\le m^{-2/5}\}$, Taylor's formula $\Lambda(\theta)=\Lambda(\theta_0)+\ii\omega(\theta-\theta_0)-\frac12\Omega(\theta-\theta_0)^2+O(m|\theta-\theta_0|^3)$
and a Gaussian integral give $T_{uv}+O(|T_{uv}|m^{-1/5})$; the integral over $J_-=-J_+$ is its complex conjugate,
as $Q_{uv}$ has real coefficients. Off $J_+\cup J_-$, $|\cos\theta-w_0|\ge c|\theta\mp\theta_0|$ because $\theta_0$
stays away from $0$ and $\pi$, and $\rho''\le-cm$ gives a modulus at most $e^{\rho(w_0)-cm^{1/5}}$.
\end{proof}

\begin{lemma}\label{lem:phase}
If $(u,v)$ and $(u+1,v)$ lie in $W_\lambda$, then $\arg T_{u+1,v}-\arg T_{uv}=\alpha_P+O(m^{-1/2})$ modulo $2\pi$.
\end{lemma}

\begin{proof}
Let $g(z)=(z-1)/X(z)$, so $Q_{u+1,v}=gQ_{uv}$. Passing from $(u,v)$ to $(u+1,v)$ changes $w_0$ by $-1/(4Pqm)$, so
$\theta_0$ moves by $O(1/m)$, and as in~\cite[Lemma~6]{Shm9} $\arg T$ moves by $\mathrm{Im}\ln g(z_0)+O(m^{-1/2})$. By
\eqref{eq:ident}, $g(z^*)=(z^*-1)\overline{X^*}/|X^*|^2=z^*/P$, so $\arg g(z^*)=\theta^*=\alpha_P$, and
$z_0-z^*=O(m^{-1/2})$.
\end{proof}

If $P$ is rational then so is $\cos\alpha_P=(2P-1)/(2P)$, and it lies in $(0,\frac12)$; by Niven's
theorem~\cite{Niv} a rational multiple of $\pi$ with rational cosine has cosine in $\{0,\pm\frac12,\pm1\}$, so
$\alpha_P/\pi$ is irrational. This is the last claim of Theorem~\ref{thm:bias}.

\section{Proof of Theorem~\ref{thm:bias}}\label{sec:proof}

The influence is Lemma~\ref{lem:sum} and the upper bound is~\eqref{eq:upper}. The lower bound is the argument
of~\cite[Section~4]{Shm9}, which we repeat with the constants of a general $P$. Fix $\varepsilon\in(0,\frac14)$,
$\lambda\ge1$ and $L\ge1$ so large that for every real $\psi_0$ at most $3\varepsilon L$ of the numbers
$\ell=0,\dots,L-1$ satisfy $|\cos(\psi_0+\ell\alpha_P)|<2\varepsilon$; this is possible because $\ell\alpha_P/\pi$ is
equidistributed modulo $1$~\cite{Weyl16}. Fix $v$ and cut the $u$ with $(u,v)\in W_\lambda$ into blocks of $L$
consecutive values. By Lemma~\ref{lem:phase}, in each block at most $3\varepsilon L$ values have
$|\cos\arg T_{uv}|<\varepsilon$ once $m$ is large. On $W_\lambda$, $\beta_{u+1}/\beta_u=(m-u)q/((u+1)P)=1+O(m^{-1/2})$,
so the weights in a block differ by a factor at most $2$, and the bad values and the left-over values carry
$\beta_u\beta_v$-weight at most $6\varepsilon+o(1)$ in all.

Let $(u,v)\in W_\lambda$ have $|\cos\arg T_{uv}|\ge\varepsilon$. By Lemma~\ref{lem:laplace},
$|\tau_{uv}|\ge\varepsilon|T_{uv}|\ge c\varepsilon m^{-1/2}e^{\rho(w_0)}$, and with $2m-u-v=n_2-1$ and
Lemma~\ref{lem:exponent},
\[
r_{uv}\ge\frac{2|\tau_{uv}|}{P^{n_2-1}}-(P-q)\ge c\varepsilon m^{-1/2}P^{-(m-1)}-1 ,
\]
so $\Psi(r_{uv})\ge\log r_{uv}^2\ge2m\log\frac1P-\log m-C_\varepsilon$. By Chebyshev's inequality $W_\lambda$ has weight
at least $(1-Pq/\lambda^2)^2$, and since $\Psi\ge0$, Lemma~\ref{lem:sum} gives
\[
\frac{E_x^P}{\Inf_x^P}\ge\Bigl(\bigl(1-\tfrac{Pq}{\lambda^2}\bigr)^2-6\varepsilon-o(1)\Bigr)\Bigl(2m\log\frac1P-\log m-C_\varepsilon\Bigr).
\]
Dividing by $\log(1/\Inf_x^P)=2m\log\frac1P$ and letting $m\to\infty$, then $\lambda\to\infty$ and $\varepsilon\to0$,
gives $\liminf\ge1$. \qed

\section{Numerical checks}\label{sec:num}

Lemma~\ref{lem:sum} reduces the ratio to $(m+1)^2$ terms whose ingredients are integers over powers of the denominator
of $P$; the code computes them exactly in integer arithmetic and evaluates each term of $\Phi$ as
$a\log(1+b/a)$ with a $\log(1+\cdot)$ routine. For $t=1$ the values agree with those of~\cite{Shm9}, and for $t=0,1,2,3$
at $k\le27$ with a rational-arithmetic implementation. The quotient of $E_x(F_{k,t})/\Inf_x(F_{k,t})$ by
$(k-1)\log_2(1/P_t)$, with $P_2=\frac{11}{16}$, $P_3=\frac{21}{32}$ and $P_4=\frac{163}{256}$, is:
\begin{center}
\begin{tabular}{rrrrrrrr}
$k$ & $41$ & $81$ & $121$ & $161$ & $201$ & $251$ & $301$\\ \hline
$t=2$ & $0.849$ & $0.899$ & $0.922$ & $0.936$ & $0.946$ & $0.954$ & $0.960$\\
$t=3$ & $0.854$ & $0.902$ & $0.928$ & $0.941$ & $0.950$ & $0.958$ & $0.964$\\
$t=4$ & $0.858$ & $0.909$ & $0.931$ & $0.944$ & $0.953$ & $0.960$ & $0.966$\\
\end{tabular}
\end{center}
For $\tilde f_k$ the formula of Theorem~\ref{thm:half} agrees with the double sum at $k=13$, $41$, $81$ and $121$
to six decimals, and the identity of Lemma~\ref{lem:kraw} was checked against direct expansion for
$m\le11$. The quotient by $k-1$ is $0.466$, $0.476$, $0.482$, $0.490$, $0.496$ and $0.499$ at $k=41$, $81$, $121$,
$301$, $1001$ and $4001$, and $E_x/\Inf_x-m+\frac12\log_2m$ is $0.7796$, $0.7677$, $0.7639$, $0.7593$, $0.7572$ and
$0.7565$ at the same values of $k$. The difference from $c_0=0.75627$, multiplied by $m$, is $0.467$, $0.459$,
$0.456$, $0.453$, $0.451$ and $0.4510$ at these $k$, and $0.4509$ at $k=8001$, in line with the error term $O(1/m)$
of Theorem~\ref{thm:half}. Extrapolating linearly in $1/m$ between consecutive values of $k$ from $1001$ to $8001$
gives $0.450842$ each time, which is $\frac5{16}\log_2e$ to six places, as in Proposition~\ref{prop:next}. After that term
is removed, the remainder multiplied by $m^2$ is $0.322$, $0.311$, $0.307$, $0.303$, $0.301$, $0.301$ and $0.301$ at
$k=41,81,121,301,1001,2001,4001$.

\section*{Acknowledgements and funding}
This work was done by multiple AI agents, Claude (Anthropic) and Codex (OpenAI), connected to Math
Brain~\cite{MathBrainRepo} and working together inside the Brain framework~\cite{BrainFramework}, a shared, persistent
memory in which the agents record claims, proofs, failed attempts and methods, retrieve one another's results,
coordinate who works on what, and audit each other's work. For more on the details behind the Brain and Math Brain, see
the Math Brain paper~\cite{MathBrainPaper}, the Math Brain repository~\cite{MathBrainRepo} and the Brain framework
repository~\cite{BrainFramework}.

The research presented in this paper was supported by the European Research Council (ERC) under the European Union's
Horizon Europe research and innovation programme (grant agreement No.~101041711), by the Simons Foundation as part of
the Collaboration on the Mathematical and Scientific Foundations of Deep Learning, by Heights Labs, by Convex Nexus
Capital, by the Israel Science Foundation (grant number 2258/19), and by the Israel Science Foundation (ISF Grant
4101/25). Computational resources were provided in part by Google Cloud through its education credits programme (credit
530550103).

\section*{Declaration of competing interest}
The author holds equity in Heights Labs and in Convex Nexus Capital and has a management role at Convex Nexus Capital.
These entities are included in the funding declaration above.

\section*{Declaration of generative AI and AI-assisted technologies}
The paper comes out of Math Brain, a project on the use of AI systems in research mathematics~\cite{BrainFramework,MathBrainRepo}. The family $F_{k,t}$
is from~\cite{Shm8} and the method for $P=\frac34$ from~\cite{Shm9}. Claude Opus 5.5 (Anthropic) found
Theorems~\ref{thm:bias} and~\ref{thm:half} and their proofs, wrote the numerical checks and prepared the manuscript.
The author directed the work and is responsible for the content of the paper.

\section*{Code and data availability}
Code and outputs: \url{https://github.com/yspennstate/fei-forced-coordinate-bias}.

The paper is part of the Math Brain project, whose public repository is
\url{https://github.com/yspennstate/math-brain-public}~\cite{MathBrainRepo}.

\appendix
\section{Notation}\label{app:notation}

Symbols are listed roughly in the order in which they first appear, and the last column gives that place. The letters $a$ and $b$ are also the arguments of $\Phi$; $n$ and $j$ are generic integers in~\eqref{eq:binom} and in Stirling's formula, $j$ and $K$ are the index and the bound of the lattice sum in Remark~\ref{rem:near}, and $\nu$ is the exponent in $[z^\nu]$.

\begin{longtable}{@{}>{\raggedright\arraybackslash}p{0.19\textwidth}>{\raggedright\arraybackslash}p{0.54\textwidth}>{\raggedright\arraybackslash}p{0.2\textwidth}@{}}
\toprule
Symbol & Meaning & First appears \\
\midrule
\endhead
$f$, $n$ & a Boolean function $\{0,1\}^n\to\{-1,1\}$, and its number of variables & Section~\ref{sec:intro} \\
$S$, $\hat f(S)$, $\chi_S$ & a set of coordinates, the Fourier coefficient of $f$ at $S$, and the character of $S$ & Section~\ref{sec:intro} \\
$U$ & the spectral sample, a random set with $\Pr[U=S]=\hat f(S)^2$ & Section~\ref{sec:intro} \\
$H(\cdot)$, $\E$, $\Pr$ & Shannon entropy in bits, expectation and probability & Section~\ref{sec:intro} \\
$C_{\mathrm{FEI}}$ & the constant of the Fourier entropy-influence conjecture $H(U)\le C_{\mathrm{FEI}}\,\E|U|$ & Section~\ref{sec:intro} \\
$x$ & the hidden coordinate, the coordinate under study; it is uniform on $\{0,1\}$ & Section~\ref{sec:intro} \\
$E_x$, $\Inf_x$ & the erasure entropy $H(\mathbf 1[x\in U]\mid U\setminus\{x\})$ and the influence $\Pr[x\in U]$ at $x$ & Section~\ref{sec:intro} \\
$\mathbf 1[\cdot]$ & the indicator of an event & Section~\ref{sec:intro} \\
$h_2$, $h_e$ & the binary entropy function in bits, and in nats & \eqref{eq:upper}; proof of Theorem~\ref{thm:half} \\
$e$ & the base of the natural logarithm & \eqref{eq:upper} \\
$k$, $m$ & an odd number of inputs, and $m=(k-1)/2$ & Section~\ref{sec:intro} \\
$y$, $y_j$ & a point of $\{0,1\}^k$, and its $j$-th coordinate & Section~\ref{sec:intro} \\
$y^0$, $y^1$ & the point with $y^0_j=1$ exactly for $j\le m$, and the point that differs from it only in $y^1_{m+1}=1$ & Section~\ref{sec:intro} \\
$\MAJ_k$, $M$ & the majority of $k$ inputs, and $\MAJ_k$ with $y^1$ removed from its true set & Section~\ref{sec:intro} \\
$\tilde f_k$ & $M(y)\vee(x\wedge[y\ge y^0])$, a function of $y$ and $x$ & Section~\ref{sec:intro} \\
$t$, $F_{k,t}$ & a positive integer, and $\tilde f_k$ with $y_j$ replaced by ``at least $t$ of $2t$'' of a fresh block for $j\le m$ and by ``at least $t+1$ of $2t$'' for $j>m$ & Section~\ref{sec:intro} \\
$\mathrm{Bin}(2t,\frac12)$, $P_t$ & a binomial variable, and $\Pr[\mathrm{Bin}(2t,\frac12)\ge t]$ & Section~\ref{sec:intro} \\
$P$, $E_x^P$, $\Inf_x^P$ & a bias, and the erasure entropy and the influence of $x$ in the orthonormal basis of the product measure with $\Pr[y_j=1]=P$ for $j\le m$ and $1-P$ for $j>m$ & Section~\ref{sec:intro} \\
$\alpha_P$ & $\arccos\frac{2P-1}{2P}$, the rotation angle & Section~\ref{sec:intro} \\
$\Psi(r)$, $r$ & $r^2\log_2(1+r^{-2})+\log_2(1+r^2)$, and its argument & Section~\ref{sec:intro} \\
$C$, $c$ & positive constants that may change from line to line: absolute in Section~\ref{sec:half}, depending on $P$ and $\lambda$ in Sections~\ref{sec:circle} and~\ref{sec:proof} & Theorem~\ref{thm:half} \\
$c_0$ & $\frac34\log_2e-\frac12\log_2\frac\pi2$ & Theorem~\ref{thm:half} \\
$b$ & an index, $0\le b\le m-1$, with $u+v=2b+1$ & Theorem~\ref{thm:half} \\
$\log$, $\ln$ & the logarithms to base $2$ and to base $e$ & Section~\ref{sec:sum} \\
$\Phi(a,b)$ & $a\log\frac{a+b}a+b\log\frac{a+b}b$, so that $\Psi(r)=\Phi(r^2,1)$ & Section~\ref{sec:sum} \\
$q$, $\sigma$ & $1-P$, and $\sqrt{Pq}$ & Section~\ref{sec:sum} \\
$\phi_j$, $\mu_j$ & the basis function $(y_j-\mu_j)/\sigma$ of the coordinate $j$, and the mean of $y_j$, $P$ for $j\le m$ and $q$ for $j>m$ & Section~\ref{sec:sum} \\
$D$, $\bar A$, $N$ & the parts of $\tilde f_k=\bar A+(2x-1)D$ in $\pm1$ values, and the indicator of $|y|\ge m+1$ & Section~\ref{sec:sum} \\
$u$, $v$ & $|S\cap\{1,\dots,m\}|$ and $|S\cap\{m+2,\dots,k\}|$ & Section~\ref{sec:sum} \\
$z$, $X(z)$, $Y(z)$ & a complex variable, $q+Pz$ and $P+qz$ & Section~\ref{sec:sum} \\
$[z^\nu]$ & the coefficient of $z^\nu$ in a polynomial & Section~\ref{sec:sum} \\
$\Pi_{uv}(z)$, $\tau_{uv}$ & $(z-1)^{u+v}X(z)^{m-u}Y(z)^{m+1-v}$, and the sum of its coefficients at $z^\nu$, $\nu\ge m+1$ & Section~\ref{sec:sum} \\
$y'$ & a point of $\{0,1\}^k$ & proof of Lemma~\ref{lem:coef} \\
$\beta_u$ & $\binom mu q^uP^{m-u}$ & Section~\ref{sec:sum} \\
$r_{uv}$ & $|\hat{\bar A}(S)|/|\hat D(S)|$ for the sets $S\not\ni m+1$ with counts $u$, $v$ & Section~\ref{sec:sum} \\
$Q_{uv}(z)$ & $\Pi_{uv}(z)/(z-1)$, a polynomial when $u+v\ge1$ & Section~\ref{sec:sum} \\
$a$, $\kappa_a$ & an index, $0\le a\le2m$, and $[z^m](1-z)^a(1+z)^{2m-a}$, a value of a Krawtchouk polynomial & Section~\ref{sec:half} \\
$\xi$ & a formal variable & proof of Lemma~\ref{lem:kraw} \\
$s$ & $u+v$; under the weights $\beta_u\beta_v$ at $P=\frac12$ it is binomial with parameters $2m$ and $\frac12$ & proof of Theorem~\ref{thm:half} \\
$p_b$, $d_b$, $\delta_b$ & $\binom{2m}{2b+1}4^{-m}=\Pr[s=2b+1]$, $2b-m$, and $d_b^2/m$ & proof of Theorem~\ref{thm:half} \\
$\zeta$ & the distance of the argument of $h_2$ or $h_e$ from $\frac12$; for the binomial coefficients, $d_b/(2m)$ & proof of Theorem~\ref{thm:half} \\
$B$ & the set of $b$ with $|d_b|\le m/4$ & proof of Theorem~\ref{thm:half} \\
$\Delta_1$, $\Delta_2$, $\Delta_3$, $\Delta_b$ & the errors in Stirling's formula for $\binom mb$, $\binom{2m}{2b}$ and $\binom{2m}m$, and the error in $\log r^2$ & proof of Theorem~\ref{thm:half} \\
$\eta$, $c_1(m)$ & $P-\frac12$, and the constant term of the first-order form of the ratio & Remark~\ref{rem:near} \\
$Z$, $\gamma$ & a standard normal variable, and Euler's constant & Remark~\ref{rem:near} \\
$\ii$, $\theta$, $\overline{X}$ & the imaginary unit, an angle with $z=e^{\ii\theta}$, and the complex conjugate of $X$ & Section~\ref{sec:circle} \\
$w$, $R(w)$ & $\cos\theta$, and $P^2+q^2+2Pqw=|X(z)|^2=|Y(z)|^2$ & Section~\ref{sec:circle} \\
$n_1$, $n_2$ & $u+v-1$ and $2m+1-u-v$ & Section~\ref{sec:circle} \\
$\rho(w)$, $w_0$ & $\ln|Q_{uv}(e^{\ii\theta})|$ as a function of $w=\cos\theta$, and its maximiser & Section~\ref{sec:circle} \\
$\lambda$, $W_\lambda$ & a parameter $\lambda\ge1$, and the set of $(u,v)$ with $|u-mq|,|v-mq|\le\lambda\sqrt m$ & Section~\ref{sec:circle} \\
$w^*$, $\theta^*$, $z^*$, $X^*$ & $\cos\alpha_P$, $\alpha_P$, $e^{\ii\theta^*}$ and $X(z^*)$ & Section~\ref{sec:circle} \\
$\theta_0$, $z_0$ & $\arccos w_0$ and $e^{\ii\theta_0}$ & Section~\ref{sec:circle} \\
$G(n_1)$ & the left side of Lemma~\ref{lem:exponent} as a function of a real $n_1$ & proof of Lemma~\ref{lem:exponent} \\
$\Lambda(\theta)$ & $\ln Q_{uv}(e^{\ii\theta})-\ii m\theta$ near $\theta^*$ & Section~\ref{sec:circle} \\
$\omega$, $\Omega$, $T_{uv}$ & $\mathrm{Im}\,\Lambda'(\theta_0)$, $-\Lambda''(\theta_0)$, and the saddle-point term $(2\pi\Omega)^{-1/2}\exp(\Lambda(\theta_0)-\omega^2/(2\Omega))$ & Section~\ref{sec:circle} \\
$J_+$, $J_-$ & the arc $|\theta-\theta_0|\le m^{-2/5}$, and its reflection $-J_+$ & proof of Lemma~\ref{lem:laplace} \\
$g(z)$ & $(z-1)/X(z)$, so that $Q_{u+1,v}=gQ_{uv}$ & proof of Lemma~\ref{lem:phase} \\
$\varepsilon$, $C_\varepsilon$ & a small parameter in $(0,\frac14)$, and a constant depending on $\varepsilon$ & Section~\ref{sec:proof} \\
$L$, $\ell$, $\psi_0$ & a block length, an index $0\le\ell\le L-1$, and a real phase & Section~\ref{sec:proof} \\
\bottomrule
\end{longtable}

\begin{thebibliography}{99}
\bibitem{ACKSW} S.~Arunachalam, S.~Chakraborty, M.~Kouck\'y, N.~Saurabh, R.~de~Wolf, Improved bounds on Fourier
entropy and min-entropy, ACM Trans. Comput. Theory 13 (2021), no.~4, 1--40, doi:10.1145/3470860; arXiv:1809.09819.
\bibitem{MI26a} Mosaic Intelligence, A certified lower bound $C_{71}\ge 6.4901128435233943$ for the Fourier entropy-influence constant, Zenodo, 2026, doi:10.5281/zenodo.20794146.
\bibitem{MI26b} Mosaic Intelligence, A certified $n=17$ lower bound for the Fourier entropy-influence constant, Zenodo, 2026, doi:10.5281/zenodo.21113935.
\bibitem{Num26} Numaro, Improved lower bound $C_{71}>6.521845710923046575$ for the Fourier entropy-influence constant ($n=18$ explicit construction), certificate archive, Zenodo, 2026, doi:10.5281/zenodo.21497769.
\bibitem{Shalev} G.~Shalev, On the Fourier entropy influence conjecture for extremal classes, arXiv:1806.03646, 2018.
\bibitem{Cha25} F.~Chang, A lower bound for the Fourier entropy of Boolean functions on the biased hypercube,
arXiv:2511.07739, 2025.
\bibitem{CKLS} S.~Chakraborty, R.~Kulkarni, S.~V.~Lokam, N.~Saurabh, Upper bounds on Fourier entropy, Theoret.
Comput. Sci. 654 (2016), 92--112.
\bibitem{FS09} P.~Flajolet, R.~Sedgewick, \emph{Analytic Combinatorics}, Cambridge University Press, 2009.
\bibitem{GMP} M.~J.~Gonz\'alez, P.~MacManus, M.~C.~Pereyra, Further evidence towards the Fourier entropy-influence
conjecture, arXiv:2606.00246, 2026.
\bibitem{FK96} E.~Friedgut, G.~Kalai, Every monotone graph property has a sharp threshold, Proc. Amer. Math. Soc.
124 (1996), 2993--3002.
\bibitem{Han} X.~Han, A new bound for the Fourier-entropy-influence conjecture, Combinatorica 45 (2025), no.~1,
Art.~4, doi:10.1007/s00493-024-00133-z; arXiv:2312.08271.
\bibitem{Hoe63} W.~Hoeffding, Probability inequalities for sums of bounded random variables, J. Amer. Statist. Assoc.
58 (1963), 13--30.
\bibitem{Hod} R.~Hod, Improved lower bounds for the Fourier entropy/influence conjecture via lexicographic functions,
arXiv:1711.00762, 2017.
\bibitem{KKLMS} E.~Kelman, G.~Kindler, N.~Lifshitz, D.~Minzer, M.~Safra, Towards a proof of the Fourier-entropy
conjecture?, Geom. Funct. Anal. 30 (2020), 1097--1138.
\bibitem{KMS12} N.~Keller, E.~Mossel, T.~Schlank, A note on the entropy/influence conjecture, Discrete Math. 312
(2012), 3364--3372, arXiv:1105.2651.
\bibitem{LH25} P.~Li, G.~Han, Strengthening Han's Fourier entropy-influence inequality via an
information-theoretic proof, arXiv:2512.03117, 2025.
\bibitem{MS77} F.~J.~MacWilliams, N.~J.~A.~Sloane, \emph{The Theory of Error-Correcting Codes}, North-Holland, 1977.
\bibitem{OD14} R.~O'Donnell, \emph{Analysis of Boolean Functions}, Cambridge University Press, 2014.
\bibitem{OT13} R.~O'Donnell, L.-Y.~Tan, A composition theorem for the Fourier entropy-influence conjecture, ICALP 2013,
Lecture Notes in Comput. Sci. 7965, 780--791, doi:10.1007/978-3-642-39206-1\_66.
\bibitem{OWZ} R.~O'Donnell, J.~Wright, Y.~Zhou, The Fourier entropy-influence conjecture for certain classes of
Boolean functions, ICALP 2011, Lecture Notes in Comput. Sci. 6755, 330--341.
\bibitem{Rob55} H.~Robbins, A remark on Stirling's formula, Amer. Math. Monthly 62 (1955), 26--29.
\bibitem{VW08} S.~Verd\'u, T.~Weissman, The information lost in erasures, IEEE Trans. Inform. Theory 54 (2008),
5030--5058.
\bibitem{WWW} A.~Wan, J.~Wright, C.~Wu, Decision trees, protocols, and the Fourier entropy-influence conjecture, ITCS 2014,
arXiv:1312.3003.
\bibitem{Weyl16} H.~Weyl, \"Uber die Gleichverteilung von Zahlen mod.~Eins, Math. Ann. 77 (1916), 313--352.
\bibitem{Niv} I.~Niven, \emph{Irrational Numbers}, Carus Mathematical Monographs 11, Mathematical Association of
America, 1956, Corollary~3.12.
\bibitem{Shm1} Y.~Shmalo, The one-coordinate step in the Fourier entropy-influence conjecture,
\url{https://github.com/yspennstate/fourier-entropy-influence-one-coordinate}, doi:\href{https://doi.org/10.5281/zenodo.24125404}{10.5281/zenodo.24125404}, 2026.
\bibitem{Shm8} Y.~Shmalo, Hidden coordinates and selection rules in the Fourier entropy-influence conjecture,
\url{https://github.com/yspennstate/fei-hidden-coordinates}, doi:\href{https://doi.org/10.5281/zenodo.24125561}{10.5281/zenodo.24125561}, 2026.
\bibitem{Shm9} Y.~Shmalo, A forced coordinate with nearly maximal erasure entropy,
\url{https://github.com/yspennstate/fei-forced-coordinate}, doi:\href{https://doi.org/10.5281/zenodo.24126356}{10.5281/zenodo.24126356}, 2026.
\bibitem{ShmLTF} Y.~Shmalo, Fourier entropy of threshold functions and their compositions,
arXiv:2610.16069, 2026.
\bibitem{MathBrainRepo} Y.~Shmalo, \emph{Math Brain: public snapshot, code and tests}, 2026,
\url{https://github.com/yspennstate/math-brain-public}.
\bibitem{MathBrainPaper} Y.~Shmalo, \emph{Math Brain: shared, auditable memory for multi-agent mathematical research},
October 2026, \url{https://github.com/yspennstate/math-brain-public/tree/main/papers/math-brain-2026-10-24}.
\bibitem{BrainFramework} Y.~Shmalo, \emph{Brain Framework}, public software repository, 2026,
\url{https://github.com/yspennstate/brain-framework-public}.
\end{thebibliography}

\end{document}
